% !TeX root = ../papers/nc-p3.tex
\section{十三次元の交叉を与える楕円曲線の非自明な二捩れ}
\label{sec:elliptic-two-torsion-crossing}
\begin{theorem}\label{thm:elliptic-two-torsion-crossing}
安定な二元障害の一般点を通る14次元成分と15次元成分の共通13次元部分は、
滑らかな楕円鉛筆 $(E,L,M,W)$ で $ML^{-1}$ が位数ちょうど2であるものの閉包と、
その点の近くで一致する。この滑らかな楕円曲線上の族は既約で、全て標準Koszul左右AS正則である。
有限二元障害の各向きの軌跡はその閉包に含まれる。特に共通部分の幾何学的一般点では
二つの反対鉛筆は同じ滑らかな楕円四次曲線を二捩れだけ異なる二束で埋め込んだものである。
\end{theorem}
\begin{proof}[四本の対角と二捩れ条件]
$E$ の原点を固定し、$\operatorname{Pic}^d(E)$ を因子の和で $E$ に同一視する。
$L,M$ に対応する点を $l,m$ とし、$U=H^0(L),Q=H^0(M)$ とする。
$K_0=\ker(U\otimes Q\to H^0(LM))$ は $E^2$ 上の対角で消える切断の空間である。
四変数を $x,x',y,y'$ とし、$\mathcal B=L\boxtimes L\boxtimes M\boxtimes M$ から
四つの交差対角 $x=y,x'=y,x=y',x'=y'$ を引いた線束を取る。
$K_0\otimes K_0$ と片側交換後の同じ空間との交わりは $H^0(\mathcal B)$
（四対角を引いた線束をここでは $\mathcal B$ と書く）である。
実際異なる素因子全てでの消滅はその和での消滅であり、二つの独立な組の対角を引くときは外積の切断空間となる。

この線束を $(x,x',y,y')=(z,z,-z,-z)$ の平行移動に制限すると次数は零である。
そのNéron--Severi形式は、$u=x+x'+y+y'$, $v=x-x'$, $w=y-y'$ として
$u^2+v^2+w^2$ である。実際四座標の各四次から四つの座標差の平方を引けばこの式となる。
従って上の楕円部分群を根基に持ち、平行移動しても制限の同型類は変わらず、
原点を通るものにおける類は $2(l-m)\in\operatorname{Pic}^0(E)$ である。
これが非零なら $H^0(\mathcal B)=0$ となる。

$\eta=m-l\ne0$ が二捩れとする。$b\in E[2]$ の各々に対し、
線束 $\mathcal B$ は次に同型である。
\[
 u^*\mathcal O_E([l-b])\otimes v^*\mathcal O_E([b])
                  \otimes w^*\mathcal O_E([b+\eta]).
\]
上の平方形式で同じNéron--Severi類を持ち、四座標軸への制限の次数2の類は
$l-b+b,l-b-b,l-b+b+\eta,l-b-b-\eta$ で、これは $l,l,m,m$ となる。
積多様体の $\operatorname{Pic}^0$ が各軸の積であることから同型が従う。
各点を切る切断を引き戻した積に四交差対角の切断を掛ける。
$b=0$ のとき $x=x'$ で一位の零、$y=y'$ の一般点で非零なので指標 $-+$ を持つ。
$b=\eta$ のときは逆の $+-$ となる。実際除数は各交換で保たれ、同一束二因子の
自然な交換作用に対する符号は固定対角での位数の偶奇である。
他の二つの $b$ は $++$ の相異なる除数を与える。
これで四つの切断は独立で全空間を張る。対角楕円部分群で割った後の
$(u,v,w)$ は行列式の絶対値2の同種写像で次数4、主偏極の積の引戻しの切断次元は4だからである。

$K=\ker(U\otimes Q\to H^0(LM)/W)$ は $K_0$ を含む。
二つの混合切断は $\bigwedge^2 K$ にあり、反対指標のSchur微分にそれぞれ非零の余核を与える。
これは任意の基点自由二次元鉛筆 $W$ で成り立つ。
非二捩れのときに消えると示したのは $K_0$ による四対角の切断であって、
個別の $W$ に由来する他の余核をこの議論で排除する必要はない。
\end{proof}
\begin{proof}[二元窓の基点自由鉛筆と二角形の平滑化]
一般二元点で $\ell$ は異なる二根 $a,b$ を持ち $c$ と互いに素である。
二本の射影直線を二節点で結んだ二角形 $E_0$ を取り、第一成分の節点を $a,b$ とする。
多重次数 $(4,0)$ の束 $L_0$ の貼合せを、主成分への制限空間が
$\ell P_2+kc\subset P_4$ になるように取る。この空間は $a,b$ での値の比を
$c(a),c(b)$ と同じにする四次式の全体である。多重次数 $(3,1)$ の $M_0$ は
主成分で全 $P_3$、$L_0M_0$ は全 $P_7$ を与える。残る成分での一次式は二端の値で一意に定まる。
写像
\[
 J:P_7\longrightarrow P_5,\qquad J(\ell h)=h,\quad J(cq)=L_Y(q)
\]
を取る。ここで $L_Y=G^{-1}H_{5,3}(Y)$ である。
二つの制限は交叉 $c\ell P_1$ で一致する。これは $Y\mathbin{\lrcorner}\ell=T\mathbin{\lrcorner}c$ による。
$J$ は全射で、積をその二次元核 $W_0$ で割ると元の二元辺になる。
$W_0$ の基点があれば、その点の評価は $J$ を経由する。従って二束の切断の積を
評価した純関係を与える。二束は全ての点で生成されるのでこの関係は非零となり、
有限Frobenius条件に反する。よって $W_0$ は二角形全体で基点自由である。

平滑化を具体的に取る。Legendre族 $y^2=x(x-1)(x-t)$ の原点の二重点を一回解消する。
全空間の二次初期式 $y^2+x^2-tx$ は非退化で、特別ファイバーの二つの接線は相異なる。
従って特別ファイバーは重複度一の主成分と例外成分が横断的に二回交わる上の二角形である。
無限遠の零切断 $O$ は主成分、一般ファイバーで位数2の切断 $P=(0,0)$ は例外成分の滑らかな点に着く。
$\mathcal L=\mathcal O(4O)$ と $\mathcal M=\mathcal L(P-O)$ は所望の多重次数で、
一般ファイバーでの差は非自明な二捩れである。
二角形上の指定された束対はこの特殊対に同型である。各束の貼合せの自由度は一つの単元であり、
両成分の二端を固定する二つの拡大縮小の作用の指数行列は、座標の向きを除き
$\left(\begin{smallmatrix}4&0\\3&1\end{smallmatrix}\right)$ となる。
従って代数閉体上で両単元を同時に合わせられる。

三束 $L_0,M_0,L_0M_0$ の一次コホモロジーは零である。双対束の切断は
正次数成分で消え、次数零の成分にも節点から零が伝わるためである。
基底変更定理で切断束と鉛筆 $W_0$ を持ち上げ、開条件により基点自由な一般鉛筆を得る。
二次辺は特殊点でも階数6を保つので族として延びる。
一般楕円鉛筆の最上テンソルの射影極限は特殊点の四次交叉にあり、その交叉は
一次元だから二元窓自身に一致する。この平滑化が二元パラメータの稠密開集合で働き、
既約性から軌跡全体の閉包含を得る。
\end{proof}
\begin{proof}[交叉成分の同定]
滑らかな楕円曲線と指定した非零二捩れ点の族は既約である。
実際指定点を $(0,0)$ とする $y^2=x(x-r)(x-s)$ の既約なパラメータで全てを覆える。
平行移動で次数4の $L$ を固定すれば $M=L\otimes\eta$ も決まり、その後の自動同型群は有限である。
$H^0(LM)$ の八次元空間の二次元鉛筆は12次元で、曲線の一変数と合わせて13次元となる。
反対鉛筆と積からのデータの復元により像も同じ次元である。
定理\ref{thm:elliptic-pencil-all-as}が全ファイバーの正則性を与える。
この軌跡は二元点の近傍で両対称であり、四対角の切断が混合接の跳躍を強制する。
前定理までの二成分模型では、そのような両対称点は13次元の共通部分に入る。
平滑化による閉包含、次元と既約性から局所的な一致を得る。
\path{computations/verify_elliptic_two_torsion_crossing.py}で四対角の形式、
解消の接錐、元の二元標本における $J$ と鉛筆を検算できる。
\end{proof}
